% Source: https://artofproblemsolving.com/community/c3h2449344_interesting_snowman_problem_that_my_math_teacher_made
% Stable identity: aops-topic-2449344-q01
A snowman is made of three snowballs stacked on top of one another. The snowballs are perfect spheres whose diameters are in the ratio $1:2:4$, and the largest, at the bottom, has diameter $36$ inches. The snowman fits exactly inside a closed glass cylinder: the radius of the cylinder equals the radius of the largest snowball, the bottom snowball rests on the base of the cylinder, each snowball rests on the one below it, and the top snowball touches the lid. The packed snow of the snowman melts into water at a ratio of $6$ to $1$ by volume (every $6$ cubic inches of snow give $1$ cubic inch of water). After the whole snowman has melted, how far below the top of the cylinder is the surface of the water?
